\documentclass[11pt]{article}

\usepackage[margin=1in]{geometry}
\usepackage{amsmath, amssymb}
\usepackage{booktabs}
\usepackage{hyperref}
\usepackage{url}
\usepackage{enumitem}
\setlist{leftmargin=*}
\setlength{\parskip}{0.35em}
\setlength{\parindent}{0pt}

\title{\textbf{Stage 5: Optional Extensions}\\
Heuristic Improvements, Application-Level Simulation, and Future Work}
\author{Team Members: Name 1, Name 2, Name 3}
\date{CS240: Algorithm Design and Analysis, Spring 2026}

\begin{document}

\maketitle

\section{Purpose of the Extension Stage}

The required project already implements and evaluates a classical graph
partitioning method. This optional stage describes extensions that go beyond
the minimum requirements. The goal is not to claim state-of-the-art graph
partitioning performance, but to show how the core algorithmic idea can be
adapted, improved, and connected more deeply to the modern application domain
of distributed optimization.

The project includes one implemented extension and several natural future
extensions. The implemented extension is a lightweight consensus-averaging
simulation that connects graph partitioning to cross-worker communication
cost. The future extensions include weighted communication costs, stronger
preprocessing, sparse spectral methods, multilevel partitioning, and more
realistic distributed optimization simulations.

\section{Implemented Extension: Consensus Communication Simulation}

The main report evaluates graph partitioning not only by pure graph metrics
but also by a small distributed consensus simulation. Each vertex holds a
scalar value, and neighboring vertices repeatedly exchange information. The
update resembles standard average consensus:
\[
    x_i^{(t+1)}
    =
    x_i^{(t)}
    +
    \alpha
    \sum_{j:(i,j)\in E}
    \left(x_j^{(t)}-x_i^{(t)}\right),
\]
where $\alpha$ is chosen based on the maximum degree. The target value is the
initial average of all node values.

The partition is used to count how many messages cross worker boundaries. For
each edge whose endpoints are assigned to different partitions, the simulation
counts two directed cross-worker messages per iteration. This produces a
communication metric that is directly related to the cut:
\[
    \text{cross messages}
    \approx
    2T \cdot |\{(u,v)\in E:p(u)\neq p(v)\}|,
\]
where $T$ is the number of consensus iterations.

This extension gives the graph partitioning objective an application-level
interpretation. Lower cut weight means fewer cross-worker messages, which is
exactly the kind of cost that distributed systems try to reduce.

\section{Observed Extension Results}

The average cross-worker message counts are:
\[
    \text{greedy}=15460,\qquad
    \text{spectral}=11360,\qquad
    \text{Kernighan--Lin}=10440.
\]
These results mirror the cut-weight comparison. Kernighan--Lin refinement has
the lowest average cut and therefore the lowest average cross-worker traffic.
This supports the idea that a classical graph partitioning algorithm can be
used as a communication-placement tool in distributed optimization.

The consensus error itself is nearly identical across partitions on the same
graph. This is expected because the simulation keeps the original graph
topology and averaging dynamics unchanged. The partition affects only where
messages are placed, not which vertices communicate. This result is useful
because it clarifies what the current extension measures: communication
placement cost rather than convergence acceleration.

\section{Possible Extension 1: Weighted Communication Costs}

The current experiments use mostly unweighted or uniformly weighted graphs. A
more realistic distributed optimization system may have nonuniform
communication costs. For example, some dependencies may require large tensor
transfers, some worker pairs may communicate through slower network links, and
some variables may be more strongly coupled than others.

An extension would assign meaningful weights $w_{uv}$ to edges and optimize
\[
    \sum_{(u,v)\in E,\;p(u)\neq p(v)} w_{uv}.
\]
Weights could be generated from:
\begin{itemize}
    \item data dependency strength;
    \item number of bytes transferred between tasks;
    \item historical communication frequency;
    \item network latency between worker groups.
\end{itemize}
This extension would make the partitioning objective more faithful to real
distributed systems.

\section{Possible Extension 2: Improved Preprocessing}

Graph preprocessing can strongly influence partition quality and runtime. A
simple preprocessing extension is to remove isolated vertices, merge duplicate
or nearly duplicate nodes, normalize edge weights, and order vertices by
degree or locality before partitioning.

For the greedy baseline, better ordering can improve quality. The current
implementation sorts vertices by weighted degree. More advanced orderings
could use breadth-first traversal, graph degeneracy ordering, or spatial
coordinates when available. For spectral methods, preprocessing could ensure
that disconnected components are handled carefully, since Laplacian eigenvectors
can be ambiguous when the graph has multiple components.

\section{Possible Extension 3: Sparse Spectral Partitioning}

The current spectral implementation uses dense eigen-decomposition, which is
simple but not scalable. A stronger version would use sparse matrices and
iterative eigensolvers to compute only the Fiedler vector. This would reduce
memory use and make the method applicable to larger sparse graphs.

The algorithmic change would be:
\begin{enumerate}
    \item build the Laplacian as a sparse matrix;
    \item use an iterative solver for the smallest nonzero eigenpairs;
    \item recursively bisect using the computed Fiedler vector;
    \item fall back to simple splitting for tiny or disconnected components.
\end{enumerate}

This extension directly addresses scalability and would make the spectral
baseline more competitive for large graph instances.

\section{Possible Extension 4: Multilevel Partitioning}

A more ambitious extension is a METIS-style multilevel algorithm. Such an
algorithm has three phases:
\begin{enumerate}
    \item \textbf{Coarsening}: repeatedly match and contract vertices to create
    smaller graphs;
    \item \textbf{Initial partitioning}: partition the smallest graph;
    \item \textbf{Uncoarsening and refinement}: project the partition back to
    larger graphs and refine it at each level.
\end{enumerate}

This would combine global scalability with local improvement. Kernighan--Lin
or a related local-search method could be used during the refinement phase.
This extension is natural because it connects the reproduced algorithm to one
of the most successful practical graph partitioning frameworks.

\section{Possible Extension 5: Stronger Distributed Optimization Simulation}

The current consensus simulation counts cross-worker messages but does not
make cross-worker communication slower or less reliable. A stronger extension
would introduce a cost model:
\begin{itemize}
    \item cross-worker messages have higher latency;
    \item cross-worker communication is allowed only under a budget;
    \item delayed messages affect the update rule;
    \item different partitions lead to different wall-clock convergence times.
\end{itemize}

For example, a delayed consensus update could distinguish local messages from
remote messages. Local messages update immediately, while remote messages are
applied after a delay. Then partitions with fewer remote edges may not only
reduce communication count but also improve convergence per unit time.

Another extension is distributed gradient descent. Each vertex could represent
a local objective $f_i(x)$, and agents could perform local gradient steps plus
neighbor averaging. The partition would then determine which synchronization
events are local and which are remote.

\section{Project-Level Takeaway}

The optional extensions strengthen the link between classical algorithms and
modern applications. The implemented consensus simulation already shows that
lower graph cuts translate into fewer cross-worker messages. The proposed
extensions would make the system model more realistic and the algorithms more
scalable.

The most valuable next step would be a weighted communication model combined
with sparse spectral partitioning. This would preserve the current project
structure while addressing two important limitations: uniform edge costs and
dense eigen-decomposition. A larger future version could then add multilevel
coarsening and compare directly with METIS-style tools.

\end{document}
